\section{Executive summary}

This report examines \textbf{liquidated damages (LD) tied to availability and
performance guarantees} for HVDC transmission systems, and how they should be
framed for American Terawatt's projects supplying large industrial and
data-center loads.

For an HVDC scheme, the commercially relevant guarantees are \emph{operational}:
the fraction of energy the link can transmit over a reporting period. The
industry standard for measuring this is CIGRE Technical Brochure~956 (2025),
which defines forced and scheduled energy unavailability (FEU, SEU) and the
forced outage rate (FOR). When the measured value exceeds the contractually
guaranteed threshold, liquidated damages become payable.

\textbf{Key points}
\begin{itemize}
  \item \textbf{Guaranteed values (realistic, OEM-backable):} FEU $\approx$
        0.5\,\%, SEU $\approx$ 1\,\%, and a bipolar forced-outage rate of order
        1 per 10 years. Very aggressive figures (e.g.\ 1 bipolar outages per 20 years) are
        \emph{commercially driven, not technically justified}, and no OEM will
        guarantee them with high penalty costs.
  \item \textbf{LD market ranges (\raise.17ex\hbox{$\scriptstyle\sim$}1\,GW project):}
        \$0.1--1.2\,M per additional forced outage; \$75--600\,k per increment
        of forced energy unavailability; overall LD cap of 3--5\,\% of contract
        value in the last years. Prior to that, LD cap of up to 20\,\% were very common.
  \item \textbf{Availability gap:} a single HVDC link cannot, on its own, meet a
        data-center Tier~III target ($\approx$99.982\,\%). Meeting it requires
        redundancy (dual full-capacity circuits), on-site back-up generation, or
        a negotiated availability target — all of which shape the LD position.
  \item \textbf{LD design principles:} LDs should be expressed as a percentage
        of OEM supply value, set high enough to \emph{incentivise the OEM to
        fix problems}, and understood as a behavioural lever rather than full
        recovery of loss.
\end{itemize}

\textbf{Recommendation.} Base the LD position on a defensible economic argument
--- the unavailability charges American Terawatt faces under its power purchase
agreement (PPA) and the cost of the battery capacity required to compensate for
forced-unavailability events --- and argue for an LD cap linked to that overall
PPA liability rather than defaulting to the standard 3--5\,\%.
